\documentclass[11pt]{article}
\usepackage[letterpaper,margin=0.9in]{geometry}
\usepackage[T1]{fontenc}
\usepackage{lmodern,microtype,amsmath,amssymb,amsthm,mathtools}
\usepackage{graphicx,booktabs,tabularx,array,caption,xcolor,enumitem,float}
\usepackage[numbers,sort&compress]{natbib}
\usepackage{hyperref}
\hypersetup{colorlinks=true,linkcolor=blue!45!black,citecolor=blue!45!black,urlcolor=blue!45!black,
 pdftitle={Multidimensional scaling of expanding paraboloids and saddles}}
\usepackage{bookmark}
\setlength{\parskip}{0.35em}
\setlength{\parindent}{0pt}
\setlist{itemsep=2pt,topsep=4pt}
\emergencystretch=1.5em
\newtheorem{proposition}{Proposition}[section]
\newtheorem{theorem}[proposition]{Theorem}
\newtheorem{corollary}[proposition]{Corollary}
\theoremstyle{remark}\newtheorem{remark}[proposition]{Remark}
\newcommand{\R}{\mathbb R}
\newcommand{\one}{\mathbf 1}
\newcommand{\norm}[1]{\left\lVert#1\right\rVert}
\newcommand{\abs}[1]{\left|#1\right|}
\newcommand{\asinh}{\operatorname{asinh}}
\newcommand{\diag}{\operatorname{diag}}
\newcommand{\Cov}{\operatorname{Cov}}
\newcommand{\Erel}{E_{\mathrm{rel}}}
\newcommand{\Eabs}{E_{\mathrm{abs}}}
\providecommand{\reportbuilddatetime}{unknown}
\InputIfFileExists{build_info.tex}{}{}
\title{Multidimensional scaling of expanding\\paraboloids and saddles\\[0.5em]
\large Smooth geodesic distances, experiments, and asymptotic geometry}
\author{\small Source: \path{~/current_projects/grip/tools/reports/geodesic-mds/report.tex}}
\date{Build \reportbuilddatetime}

% The figure atlas preserves each existing vector figure at its native size.
% Larger pages keep the original labels readable instead of shrinking a
% twelve-panel plate into a letter-sized text column.
\newsavebox{\platebox}
\newlength{\normalpaperwidth}\setlength{\normalpaperwidth}{\paperwidth}
\newlength{\normalpaperheight}\setlength{\normalpaperheight}{\paperheight}
\newlength{\platepaperwidth}\newlength{\platepaperheight}
\newcommand{\plate}[3]{%
 \clearpage
 \sbox{\platebox}{\includegraphics{figures/#1}}%
 \setlength{\platepaperwidth}{\dimexpr\wd\platebox+1.4in\relax}%
 \setlength{\platepaperheight}{\dimexpr\ht\platebox+3in\relax}%
 \newgeometry{layoutwidth=\platepaperwidth,layoutheight=\platepaperheight,
 left=.7in,right=.7in,top=.6in,bottom=.7in}
 \paperwidth=\platepaperwidth\paperheight=\platepaperheight
 \pdfpagewidth=\paperwidth\pdfpageheight=\paperheight
 \begin{center}\usebox{\platebox}\par\vspace{.18in}
 \begin{minipage}{\dimexpr\normalpaperwidth-1.8in\relax}
 \captionsetup{font=normalsize,labelfont=bf,hypcap=false}
 \captionof{figure}{#3}\label{#2}
 \end{minipage}\end{center}
 \clearpage
 \paperwidth=\normalpaperwidth\paperheight=\normalpaperheight
 \pdfpagewidth=\paperwidth\pdfpageheight=\paperheight
 \restoregeometry
}

\begin{document}
\maketitle
\begin{abstract}
Does metric multidimensional scaling make an expanding curved surface flat?
The answer depends on the input distances, the objective, the output dimension,
and what is meant by ``flat.'' We study the graphs $z=x^2+y^2$ and $z=x^2-y^2$
over disks of radius $r$, using shortest distances along the smooth surfaces.
For the paraboloid, the distance matrix divided by $r^2$ converges to a line
metric. Both classical multidimensional scaling and global raw-stress
minimization inherit this normalized line limit. Boundary-antipodal separations
survive on the smaller scale $r$: their geodesic distance is
$\pi r-\pi^3/(96r)+O(r^{-3})$, while their fitted separation depends on the
MDS objective and sampling regime. For the saddle, the scaled geodesic matrix
converges to a four-branch tree metric. Classical three-dimensional MDS has a
provably nonplanar limit for generic samples covering all four branches, and
raw-stress computations show the same qualitative behavior. Experiments use
240 coupled samples, ten radii, two sampling measures, and smooth geodesic
solvers checked independently. All fifteen figures from the investigation are
included, together with the ambient-distance comparison, numerical estimates,
proofs, and an explanation of why an equal-axis plot of the original saddle can
look like an interval.
\end{abstract}
\vfill
\textbf{How to read the report.} Sections~\ref{sec:question}--\ref{sec:discussion}
give the definitions, results, and principal proofs. The appendices supply
geodesic formulas, the ambient-distance comparison, and numerical validation.
The complete figure atlas uses larger pages where needed to preserve readable
labels; each figure is linked from the relevant discussion.
\clearpage
{\small\setlength{\parskip}{0pt}\tableofcontents}
\clearpage

\section{The question and the main distinctions}\label{sec:question}

The motivating intuition is that MDS tries especially hard to preserve large
distances and might therefore make a curved surface progressively flatter as
the sampled domain grows. There is a useful part to this intuition: equal
\emph{relative} errors on longer targets contribute more raw stress. But it
does not imply a general preference for flat embeddings. The geometry of the
entire target distance matrix determines what a good embedding looks like.

We consider two families of compact surface patches:
\begin{align}
 \mathcal P_r&=\{(x,y,x^2+y^2):x^2+y^2\le r^2\},\\
 \mathcal S_r&=\{(x,y,x^2-y^2):x^2+y^2\le r^2\}.
\end{align}
The coefficient of each quadratic is fixed at one in the chosen coordinates.
The parameter $r$ is the radius of the \emph{base disk}, not the radius of an
intrinsic geodesic ball. Throughout the main experiments, an input
dissimilarity is a shortest path length along the smooth surface. Ambient
Euclidean chord distances appear only in the explicitly identified comparison
in Appendix~\ref{app:ambient}.

Three meanings of flattening must be separated. First, a quadratic fit in the
original units may have a smaller coefficient than the original surface.
Second, an embedding divided by its overall size may approach a plane or a
line. Third, the Gaussian curvature of the original surface may approach zero
at points far from its vertex. These statements describe different quantities
and need not agree.

The experiments and proofs give the following answers.
\begin{itemize}
 \item \textbf{Paraboloid, geodesic input.} Finite-radius three-dimensional
 fits are approximately shallower paraboloids, but their fitted quadratic
 coefficients do not trend to zero over the tested radii. After uniform
 division of all coordinates by $r^2$, the rigorous limit is a line.
 \item \textbf{Saddle, geodesic input.} Large-radius embeddings develop four
 thin branches spanning three dimensions. The classical-MDS nonplanarity
 statement is proved under a sample rank condition. The raw-stress statement
 is supported by the numerical fits, without a global optimality certificate.
 \item \textbf{Original surfaces.} Both remain smooth two-dimensional surfaces
 at every finite $r$. A narrow equal-axis plot does not mean the surface has
 ceased to extend over its full base disk.
\end{itemize}

The sampling statements also need a finite-radius interpretation. There is no
uniform probability distribution on either entire infinite-area surface. We
study a family of uniform distributions on finite patches, couple their
samples across $r$, and then take a scaled limit.

\section{Distances, objectives, and diagnostics}\label{sec:definitions}

\subsection{Raw stress and the role of a long target distance}
For $n$ sampled points let $\Delta_{ij}(r)$ be the prescribed dissimilarity,
and let $Y_i\in\R^k$ be the fitted coordinates. Fixed-target, unweighted raw
metric stress is
\begin{equation}\label{eq:stress}
 S_r(Y)=\sum_{i<j}\bigl(d_{ij}(Y)-\Delta_{ij}(r)\bigr)^2,
 \qquad d_{ij}(Y)=\norm{Y_i-Y_j}.
\end{equation}
This is the usual weighted stress formulation with every pair weight equal
to one and with the disparities fixed at the input distances
\citep[Section 2, Eq.~(1)]{mair2022}. No monotone or other transformation of the
targets is fitted.

For a single pair, write $d=d_{ij}$ and $\delta=\Delta_{ij}$. Then
\begin{equation}
 \frac{\partial(d-\delta)^2}{\partial d}=2(d-\delta),\qquad
 \frac{\partial^2(d-\delta)^2}{\partial d^2}=2.
\end{equation}
An absolute error of one contributes one, whether the target is ten or one
thousand. If the relative error is $\varepsilon=(d-\delta)/\delta$, however,
the contribution is $\delta^2\varepsilon^2$. Thus long targets receive a larger
penalty for the \emph{same relative error}, but not for the same absolute error.
They do not have a larger per-pair second derivative with respect to distance.

Uniform scaling introduces no additional shape preference. Replacing all
targets by $c\Delta$ and all coordinates by $cY$ multiplies stress by $c^2$.
Its minimizers are therefore carried to the rescaled minimizers. Any change
of normalized shape must come from a change in the relative pattern of the
target distances, rather than their common magnitude alone.

\subsection{Classical MDS is a different optimization problem}
Let $J=I-\one\one^T/n$ be the centering matrix. Classical MDS forms
\begin{equation}\label{eq:classical}
 B=-\frac12J\Delta^{\circ2}J,
\end{equation}
where $\Delta^{\circ2}$ denotes the entrywise square. If
$\lambda_1\ge\lambda_2\ge\cdots$ are its eigenvalues, the $k$-dimensional
coordinates retain the largest $k$ positive eigenvalues and multiply each
corresponding unit eigenvector by $\sqrt{\lambda_a}$. This minimizes
\begin{equation}
 \norm{B-YY^T}_F^2
\end{equation}
over centered configurations of rank at most $k$. The squared Frobenius norm
is the sum of squared matrix entries. This Gram-matrix objective,
often called strain, generally differs from raw distance stress.

The official R documentation identifies \texttt{cmdscale} as classical MDS
and describes the principal-component axes and rigid-motion ambiguity
\citep{rcmdscale}. The repository's \texttt{grip::metric.mds()} calls that
classical calculation on its prepared graph-distance matrix. The present
experiments use the corresponding spectral calculation directly on smooth
geodesic matrices, and separately fit raw stress. They do not use graph
shortest paths as a substitute for smooth geodesics.

\subsection{Sampling and numerical design}\label{sec:design}
Each radius sweep uses $n=240$ and seed 20260905 at
\begin{equation}
 r\in\{0.1,0.25,0.5,1,2,4,8,16,32,64\}.
\end{equation}
The same independent uniform radial quantiles $U_i$ and angles
$\Theta_i\sim\operatorname{Unif}[0,2\pi)$ are reused across radii and surfaces.
For a realized quantile $\xi_i$, uniform sampling in the base disk uses
$\rho_i=r\sqrt{\xi_i}$. Both surfaces have area density
$\rho\sqrt{1+4\rho^2}\,d\rho\,d\theta$. Uniform surface-area sampling therefore uses
\begin{equation}\label{eq:area-sampling}
 \rho_i=\frac12\sqrt{\left[1+\xi_i\{(1+4r^2)^{3/2}-1\}\right]^{2/3}-1}.
\end{equation}
In the second family, $\rho_i/r\to \xi_i^{1/3}$ and
$\rho_i^2/r^2\to \xi_i^{2/3}$. The limiting distribution differs from that of
uniform base-disk sampling, but the geometric bounds below apply to both.

Classical MDS is computed in two and three dimensions. For each raw-stress fit,
six starts are optimized and the lowest stress is retained. The two-dimensional
starts are classical MDS, the three coordinate-plane projections of the
original sample, and two random configurations. The three-dimensional starts
are classical MDS, the original sample, two two-dimensional fits with small
random third-coordinate perturbations, and two random configurations.
The perturbations avoid leaving a three-dimensional solver in an exactly
planar stationary subspace. Each start's overall scale is initially profiled
by least squares.

Optimization uses L-BFGS-B with an analytic gradient, maximum 4000 iterations,
and objective and gradient tolerances $10^{-14}$ and $10^{-10}$. Targets and
coordinates are divided by the RMS input pair distance during fitting, and
the optimized function is mean squared pair residual. These common factors do
not change the minimizing shape. Multiple starts address, but do not eliminate,
the local-minimum issue discussed by \citet[Section 3.1]{mair2022}.

Both geodesic sweeps attempted 240 optimization runs each: 20 cases, two output
dimensions, six starts. All 480 runs reported success and all 80 selected fits
were retained. The ambient comparison attempted 120 two-dimensional runs;
all succeeded and 20 fits were selected. Its three-dimensional solution is
known exactly and did not require numerical stress minimization. No failed
fits were omitted from the displayed results. Successful termination does
not prove global optimality.

\subsection{Distance error, dimensionality, and shape coefficients}
With $M=n(n-1)/2$, the relative and absolute root mean squared distance errors are
\begin{equation}\label{eq:errors}
 \Erel=\sqrt{\frac{\sum_{i<j}(d_{ij}-\Delta_{ij})^2}
                         {\sum_{i<j}\Delta_{ij}^2}},\qquad
 \Eabs=\sqrt{\frac1M\sum_{i<j}(d_{ij}-\Delta_{ij})^2}.
\end{equation}
The first is dimensionless; the second has the units of input distance. A
decrease in relative error does not require a decrease in absolute error.

For centered embedding coordinates let $s_1\ge s_2\ge s_3\ge0$ be their
singular values. The first principal component's variance fraction is
$s_1^2/\sum_a s_a^2$. The ratios $s_2/s_1$ and $s_3/s_2$ describe shape without
depending on the overall scale. A line limit has $s_2/s_1\to0$. A nondegenerate
plane limit has $s_3/s_2\to0$ while $s_2/s_1$ stays bounded away from zero.
The second ratio alone is insufficient to diagnose a line: two small axes
can remain comparable while both disappear relative to the first.

Quadratic shape diagnostics first center and rigidly align the embedding to
the original coordinates, without rescaling. For the paraboloid we regress
\begin{equation}\label{eq:parab-fit}
 z=a(x^2+y^2)+bx+cy+d.
\end{equation}
For the saddle we use the more general fit
\begin{equation}\label{eq:saddle-fit}
 z=ax^2+bxy+cy^2+dx+ey+f.
\end{equation}
The reported saddle coefficient norm is $\sqrt{a^2+c^2+b^2/2}$, the Frobenius
norm of the symmetric quadratic coefficient matrix. The original norm is
$\sqrt2$. The coefficient of determination is
$R^2=1-\sum_i(z_i-\widehat z_i)^2/\sum_i(z_i-\bar z)^2$.
These are descriptive regressions on finite point clouds, not estimates of
intrinsic Gaussian curvature or assertions of exact quadratic surfaces.

\section{The paraboloid: finite-radius results and the line limit}\label{sec:paraboloid}

\subsection{What the experiment shows}
Figures~\ref{fig:p3d}--\ref{fig:pshape} show the complete geodesic experiment.
Small patches give nearly planar fits, intermediate patches have three
substantial directions, and large patches become strongly elongated. The
following table gives the entire raw-stress three-dimensional base-disk sweep;
all values are outcomes of one coupled sample, not averages or confidence
intervals from independent replicates.
\input{tables/paraboloid_sweep.tex}

At $r=64$, the first principal component accounts for about $99.760\%$ of the
variance. Surface-area sampling gives the same qualitative pattern, with
$\Erel=0.001931$, $s_2/s_1=0.045085$, and $s_3/s_2=0.857030$ at that radius.
The two smaller directions remain comparable, but both become small relative
to the leading direction. This is a line limit, not a nondegenerate plane.

The quadratic diagnostic addresses the separate question of whether the
output is a shallower paraboloid in its original units.
\input{tables/paraboloid_coefficients.tex}
The original coefficient is $a=1$. Both methods produce substantially smaller
coefficients, but the coefficients do not decrease toward zero over this
sweep. At $r=64$, $R^2$ is approximately 0.9883 for classical MDS and 0.9657 for
raw stress. The leading line-limit theorem does not determine the limiting
quadratic coefficient: that coefficient depends on smaller-scale coordinates
which disappear under division by $r^2$.

\subsection{A uniform bound on all geodesic distances}
Fix unit-disk points $u_i\in\R^2$, put $q_i=\norm{u_i}^2$, and write
\begin{equation}
 P_i(r)=(ru_i,r^2q_i).
\end{equation}
Let $g_{ij}(r)$ be their smooth geodesic distance.

\begin{proposition}[Paraboloid distance bound]\label{prop:parab-bound}
For every pair and every $r>0$,
\begin{equation}\label{eq:parab-bound}
 r^2|q_i-q_j|\le g_{ij}(r)
 \le r^2|q_i-q_j|+(1+\pi)r.
\end{equation}
\end{proposition}
\begin{proof}
The length of every surface path is at least its absolute vertical displacement,
which gives the lower bound. For the upper bound, start at the endpoint with
larger base radius, move radially to the smaller radius, and then follow the
shorter circular arc to the second endpoint. Radial arclength is
$\sqrt{1+4\rho^2}\,d\rho$. The inequality
$\sqrt{1+4\rho^2}\le1+2\rho$ bounds the radial part by the height difference
plus at most $r$. The shorter circular arc has length at most $\pi r$.
The entire constructed path lies inside the base disk.
\end{proof}

The bound implies uniform convergence over pairs:
\begin{equation}\label{eq:parab-limit}
 \frac{g_{ij}(r)}{r^2}\longrightarrow |q_i-q_j|.
\end{equation}
The limit is the distance matrix of one-dimensional coordinates $q_i$.
Equal-radius points lose their angular separation on this scale. This does
not say that their unscaled geodesic distance tends to zero.

\subsection{Deriving the MDS limit directly from the objectives}
\begin{theorem}[Normalized line limit]\label{thm:line}
For a fixed finite sample with at least two distinct $q_i$, every global
raw-stress minimizer in any fixed dimension $k\ge1$ approaches the height-line
distance matrix after scaling coordinates by $r^2$. The centered Gram matrices
converge to $q_cq_c^T$, where $q_c=Jq$. The same conclusion holds for classical
MDS and for approximate raw-stress minimizers whose excess stress divided by
$r^4$ tends to zero.
\end{theorem}
\begin{proof}
Write $Y=r^2V$. Then
\begin{equation}\label{eq:rescaled-stress}
 r^{-4}S_r(r^2V)=\sum_{i<j}
 \left(\norm{V_i-V_j}-\frac{g_{ij}(r)}{r^2}\right)^2.
\end{equation}
The line candidate $V_i=(q_i,0,\ldots,0)$ has normalized stress at most
$M(1+\pi)^2/r^2$ by Proposition~\ref{prop:parab-bound}. A global minimizer does
at least as well. Its sum of squared normalized residuals therefore tends to
zero, so every individual residual tends to zero. Combined with
\eqref{eq:parab-limit}, its fitted pair distances converge to $|q_i-q_j|$.
Double-centering squared distances gives convergence of its centered Gram
matrix to $q_cq_c^T$. This matrix has rank one under the stated condition.
Allowing excess normalized stress $o(1)$ does not change the argument.

For classical MDS, the scaled input matrix $B(r)/r^4$ converges directly to
$q_cq_c^T$. Retaining its largest positive eigenvalues gives the same rank-one
limit. The principal eigenvalue is separated from zero, while every other
retained eigenvalue tends to zero on this scale.
\end{proof}

The theorem is about normalized geometry. It does not imply bounded unscaled
coordinates, monotonically decreasing finite-radius stress, a flat-disk limit,
or a limiting value of the fitted coefficient in \eqref{eq:parab-fit}.
For the coupled surface-area samples, $q_i(r)\to \xi_i^{2/3}$; the same proof
applies with the moving height-line candidate and its limiting coordinates.

\section{Boundary antipodes on the paraboloid}\label{sec:antipodal}

The line limit erases the angular separation of equal-height points. To recover
what happens to such a pair, we must examine the smaller scale $r$. Set
\begin{equation}
 p_+(r)=(r,0,r^2),\qquad p_-(r)=(-r,0,r^2),
\end{equation}
and let $g_r$ be their smooth geodesic distance. Their ambient chord distance
is $2r$. Their MDS separation is a third quantity, determined by the full
dataset, target dimension, and objective.

\subsection{The target geodesic: leading term and correction}
The boundary semicircle gives $g_r\le\pi r$. If a shortest path reaches minimum
base radius $m$, its radial travel costs at least $2(r^2-m^2)$, and its angular
travel costs at least $\pi m$. Comparing the first of these bounds to
$g_r\le\pi r$ gives $m/r\to1$. Hence
$\pi m\le g_r\le\pi r$ proves $g_r/r\to\pi$.

The exact geodesic primitives in Appendix~\ref{app:parab-geodesics} yield a
sharper result:
\begin{proposition}[Antipodal geodesic expansion]\label{prop:antipodal}
As $r\to\infty$,
\begin{equation}\label{eq:antipodal-expansion}
 g_r=\pi r-\frac{\pi^3}{96r}+O(r^{-3}),\qquad
 c=r-\frac{\pi^2}{32r}+O(r^{-3}),
\end{equation}
where $c$ is the minimum base radius of the minimizing path.
\end{proposition}
\begin{proof}
Put $t=\sqrt{r^2-c^2}$. The symmetric turning geodesic has half-angle
$H(c,t)=\pi/2$ and length $2L(c,t)$, where $H,L$ are defined in
\eqref{eq:primitives}. With $c=\sqrt{r^2-t^2}$ their expansions are
\begin{align}
 H(c,t)&=2t+\frac{-t^3/3+t/4}{r^2}+O(r^{-4}),\\
 2L(c,t)&=4rt+\frac{-4t^3/3+t/2}{r}+O(r^{-3}).
\end{align}
Writing $t_0=\pi/4$, the half-angle equation gives
$t=t_0+(t_0^3/6-t_0/8)/r^2+O(r^{-4})$. Substituting into the length and the
expression for $c$ gives \eqref{eq:antipodal-expansion}.
\end{proof}
The shortest route stays near the boundary, dipping inward by order $1/r$.
A through-vertex path has length of order $r^2$ and is much longer.

\subsection{Raw-stress MDS: a fixed-sample theorem}
\begin{theorem}[Antipodal separation under raw stress]\label{thm:raw-antipodal}
Take a fixed finite set of nonzero interior unit-disk points with distinct
squared radii, and adjoin the two boundary antipodes. Use all pair distances
with equal weights and an output dimension $k\ge2$. For a global raw-stress
minimizer,
\begin{equation}\label{eq:raw-antipodal}
 d_{\mathrm{MDS}}(p_+,p_-)=\pi r+O(1).
\end{equation}
\end{theorem}
\begin{proof}
First consider endpoints with unequal squared radii. Let their physical base
radii be $a<b$, let $h=\log(b/a)>0$, and let $\alpha\in[0,\pi]$ be their
smaller angular separation. The path
\begin{equation}
 \theta(\rho)=\theta(a)+\frac{\alpha\log(\rho/a)}h
\end{equation}
has length at most
\begin{equation}\label{eq:log-path}
 b^2-a^2+\frac h4+\frac{\alpha^2}{4h}.
\end{equation}
Indeed, its integrand is $\sqrt{4\rho^2+1+\alpha^2/h^2}$, and
$\sqrt{4\rho^2+C}\le2\rho+C/(4\rho)$ gives this bound. The ratio $b/a$ is
fixed as $r$ increases, so each unequal-radius target is
$r^2|q_i-q_j|+O(1)$.

Construct a feasible two-dimensional embedding by placing every interior
point at $(0,r^2q_i)$ and the two boundary points at
$(g_r/2,r^2)$ and $(-g_r/2,r^2)$. It fits the boundary pair exactly. Every
other pair has a nonzero fixed height-level difference, and its fitted
distance is also $r^2|q_i-q_j|+O(1)$: the transverse displacement is at most
order $r$, whose squared contribution produces only an order-one correction
to a distance of order $r^2$. The total stress of this candidate is $O(1)$
for a fixed finite sample. Thus the global minimum is $O(1)$, and in particular
\begin{equation}
 |d_{\mathrm{MDS}}(p_+,p_-)-g_r|\le\sqrt{S_r^{\min}}=O(1).
\end{equation}
Proposition~\ref{prop:antipodal} completes the proof. The construction embeds
unchanged into any $k\ge2$.
\end{proof}

The assumptions hold almost surely for a finite uniform interior sample, but
the constants can become large when two radii are very close. The theorem is
not uniform in sample size. Adding a complete boundary ring introduces many
equal-radius constraints and is outside this argument. It also does not prove
that a particular local numerical optimizer attains the required bounded
stress. The order of the fixed-sample, large-radius, and population limits
must not be silently interchanged.

\subsection{Classical MDS: the next-order finite-sample matrix}
Classical MDS does not minimize the stress used in Theorem~\ref{thm:raw-antipodal},
and its antipodal coefficient is different. For positive squared radii define
the logarithmic mean
\begin{equation}\label{eq:logmean}
 \ell(q,p)=\frac{q-p}{\log q-\log p},\qquad \ell(q,q)=q.
\end{equation}
Let $\alpha_{ij}$ denote the smaller angular separation. Expanding the smooth
geodesic primitives for fixed endpoints gives
\begin{align}
 g_{ij}(r)^2&=r^4(q_i-q_j)^2+r^2C_{ij}+O(1),\label{eq:second-order}\\
 C_{ij}&=\tfrac14(q_i-q_j)(\log q_i-\log q_j)
          +\ell(q_i,q_j)\alpha_{ij}^2.\label{eq:c-matrix}
\end{align}
For unequal radii, the first correction to the distance itself is
\begin{equation}
 g_{ij}(r)=r^2|q_i-q_j|
 +\frac{|\log(q_i/q_j)|}{8}
 +\frac{\alpha_{ij}^2}{2|\log(q_i/q_j)|}+O(r^{-2}).
\end{equation}
For equal radii, the fixed-angle turning expansion instead gives
$g_{ij}^2=r^2q_i\alpha_{ij}^2+O(1)$, agreeing with the continuous expression
for $C_{ij}$. These are fixed-endpoint expansions; the remainders need not be
uniform as unequal radii approach equality or zero.

Let $q_c=Jq$ and project off the constant and height directions:
\begin{equation}
 H=I-\frac{\one\one^T}{n}-\frac{q_cq_c^T}{\norm{q_c}^2},\qquad
 T_{\perp}=-\frac12HCH.
\end{equation}
The leading classical axis is of size $r^2$ and follows $q_c$. The smaller
axes, of size $r$, are determined by the positive eigenvalues of $T_{\perp}$.
If $(\lambda_a,v_a)$, $a=1,2$, are its leading two positive eigenpairs and the
retained subspace is separated from the discarded spectrum, then classical
three-dimensional MDS satisfies
\begin{equation}\label{eq:finite-kappa}
 \frac{d_{\mathrm{classical}}(p_+,p_-)}r\longrightarrow
 \kappa_n=\left[\sum_{a=1}^2\lambda_a
       \{v_a(p_+)-v_a(p_-)\}^2\right]^{1/2}.
\end{equation}
The two points have identical $q$, so their leading height-axis difference
vanishes. The correction to that axis contributes only $O(1)$ to their
coordinate difference and therefore disappears after division by $r$.

Using the same 240 base-disk interior points and then adjoining the pair gives
$\kappa_n=2.92190551$. This dataset has 242 points; it is distinct from the
240-point radius sweep. Direct classical fits confirm the next-order matrix.
\input{tables/antipodal_finite.tex}
At $r=1024$, the fitted ratio differs from $\kappa_n$ by less than $10^{-6}$.
The geodesic correction also agrees with $\pi^3/96\simeq0.322982$ when
$(\pi r-g_r)r$ is evaluated at large $r$.

\subsection{The rotationally symmetric population operator}
In a uniform base-disk population, $q$ is uniform on $[0,1]$ and angle is
independent and uniform. The Fourier expansion of squared circular separation is
\begin{equation}
 \alpha^2=\frac{\pi^2}{3}
       +4\sum_{m\ge1}\frac{(-1)^m}{m^2}\cos(m\alpha).
\end{equation}
In the second-order operator corresponding to \eqref{eq:c-matrix}, the two
leading transverse modes are the sine and cosine modes with $m=1$. Their
radial component solves
\begin{equation}\label{eq:operator}
 \lambda f(q)=\int_0^1\ell(q,p)f(p)\,d\mu(p),\qquad
 \int_0^1 f(q)^2\,d\mu(q)=1,
\end{equation}
with $d\mu(q)=dq$ for base-disk sampling. The embedding radius in these two
coordinates is $r\sqrt{2\lambda}\,|f(q)|$. Its boundary-antipodal coefficient is
\begin{equation}\label{eq:population-kappa}
 \kappa_\mu=2\sqrt{2\lambda}\,|f(1)|.
\end{equation}
For expanding surface-area samples the limiting measure is
$d\mu(q)=\tfrac32\sqrt q\,dq$.

Gauss--Legendre quadrature with 64, 128, 256, and 512 nodes gives the following
operator estimates. The endpoint value is obtained from the eigen-equation,
$f(1)=\lambda^{-1}\int\ell(1,p)f(p)\,d\mu(p)$, rather than by treating a
quadrature node as the boundary.
\input{tables/antipodal_population.tex}
The last two base-disk estimates differ by less than $6\times10^{-9}$.
At 256 nodes the leading $m=1$ eigenvalue is approximately 0.511361. The largest
competing radial, even-angular, and remaining odd-angular eigenvalues are
approximately $0.0000194$, $0.002891$, and $0.056818$, respectively, so none is
close to displacing the retained pair in this numerical calculation.

These constants are numerical evaluations of the population second-order
operator, not exact closed forms or rigorous quadrature error bounds. The
finite-sample spectral result above is separate. In particular, the present
work does not establish a uniform remainder theorem allowing arbitrary joint
limits $n,r\to\infty$. The values $\pi$, $2.92190551$, and $2.86038105$ answer
different, explicitly specified questions; they are not contradictory estimates
of one universal MDS distance.

\section{The saddle: four branches instead of a line}\label{sec:saddle}

\subsection{Curvature and the computed embeddings}
For a graph $z=f(x,y)$, Gaussian curvature is
\begin{equation}
 K=\frac{f_{xx}f_{yy}-f_{xy}^2}{(1+f_x^2+f_y^2)^2}.
\end{equation}
Substitution gives
\begin{equation}\label{eq:curvatures}
 K_{\mathcal P}(\rho)=\frac4{(1+4\rho^2)^2},\qquad
 K_{\mathcal S}(\rho)=-\frac4{(1+4\rho^2)^2}.
\end{equation}
Their magnitudes decay identically. The saddle is negatively curved everywhere,
but does not approach a constant negative curvature. A description such as
``asymptotically hyperbolic'' needs a specified meaning; it cannot be inferred
in the constant-curvature sense from this surface equation. The following
limit concerns base-disk expansion with distances divided by $r^2$.

Figures~\ref{fig:s3d}--\ref{fig:scompare} show that the three-dimensional MDS
outputs develop four narrow branches. Their three singular values remain
comparable. The raw-stress results for uniform base-disk sampling are:
\input{tables/saddle_sweep.tex}
At $r=64$, classical MDS gives $s_2/s_1=0.9175$ and $s_3/s_2=0.9103$.
Raw stress under surface-area sampling gives 0.9220 and 0.9177, with relative
distance error 0.061128. These values provide no evidence of collapse to a
plane or a line in the ambient embedding space. In contrast, the paraboloid
raw-stress value $s_2/s_1=0.03743$ indicates strong elongation.

The saddle quadratic coefficient norm decreases with radius
(Figure~\ref{fig:sshape}). That observation describes a dimensional fit, not
planarity. If all coordinates of a configuration are multiplied by $c$, a
quadratic coefficient in a graph regression is divided by $c$. Thus a
coefficient can tend to zero while all three axes remain comparable.
The approximate four-branch shape also need not remain an exact smooth saddle.

\subsection{Smooth shortest paths stay inside the base disk}
The entire saddle graph is complete, simply connected, and negatively curved.
Completeness follows because the induced metric dominates the base Euclidean
metric, and the graph is smooth over all of $\R^2$. The Hadamard--Cartan theorem
therefore gives a unique minimizing geodesic between every pair
\citep[Theorem 6.5.2, Corollary 6.5.3, pp.~124--125]{gorodski2022}.

For an affine geodesic parameter, write $v=\dot x$, $w=\dot y$, and
$\rho^2=x^2+y^2$. The geodesic equations are
\begin{equation}\label{eq:saddle-ode}
 \ddot x=-\frac{4x(v^2-w^2)}{1+4\rho^2},\qquad
 \ddot y=\frac{4y(v^2-w^2)}{1+4\rho^2}.
\end{equation}
Consequently,
\begin{align}
 (\rho^2)''
 &=2(v^2+w^2)-\frac{8(x^2-y^2)(v^2-w^2)}{1+4\rho^2}\\
 &\ge\frac{2(v^2+w^2)}{1+4\rho^2}>0
\end{align}
for a nonconstant geodesic. A convex function cannot exceed its largest
endpoint value on an interval. The full-surface geodesic therefore stays
inside the largest endpoint base radius and is also the shortest path on the
disk-restricted patch. This justifies computing the unrestricted smooth
geodesic for the experiment's restricted domain.

\subsection{The limiting four-branch tree metric}
Fix unit-disk points $u_i=(a_i,b_i)$ and define
\begin{equation}
 P_i(r)=(ra_i,rb_i,r^2h_i),\qquad h_i=a_i^2-b_i^2,\qquad t_i=|h_i|.
\end{equation}
There are four branches: $h>0$ with $a>0$, $h>0$ with $a<0$, $h<0$ with
$b>0$, and $h<0$ with $b<0$. They are the upper $+x$, upper $-x$, lower $+y$,
and lower $-y$ branches used for colors in the figures. Zero-height points
map to their common root. Let
\begin{equation}\label{eq:tree}
 T_{ij}=\begin{cases}
 |t_i-t_j|,&\text{points on the same branch},\\
 t_i+t_j,&\text{points on different branches}.
 \end{cases}
\end{equation}
This is the path metric of four intervals attached at one root.

\begin{proposition}[Saddle distance bound]\label{prop:tree-bound}
For smooth saddle geodesic distances,
\begin{equation}\label{eq:tree-bound}
 r^2T_{ij}\le\Delta_{ij}(r)\le r^2T_{ij}+4r.
\end{equation}
In particular, $\Delta(r)/r^2\to T$ uniformly over pairs and unit-disk endpoints.
\end{proposition}
\begin{proof}
Path length dominates total vertical variation. For two points on the same
branch this gives the height difference. Different signs of height require
crossing zero. Opposite positive-height branches require $x$ to cross zero,
where height is nonpositive; opposite negative-height branches require $y$
to cross zero, where height is nonnegative. In every different-branch case,
the required total vertical variation is at least the sum of absolute heights.

For the upper bound, move each endpoint at constant height along its hyperbola
to the appropriate coordinate semiaxis. On a positive-height branch,
$|dx/dy|\le1$ and $|y|\le r/\sqrt2$, so this horizontal connector has length
at most $r$. The negative-height argument exchanges $x$ and $y$. Next connect
the two semiaxis points monotonically on one branch or through the origin if
the branches differ. The inequality $\sqrt{1+4x^2}\le1+2|x|$ bounds this part
by the required vertical variation plus at most $2r$. Including both connectors
gives the upper bound. All segments remain inside the base disk. Zero-height
points use their straight ruling to the origin.
\end{proof}

At $r=64$, replacing the entire smooth matrix by $r^2T$ has relative distance
RMSE 0.001588 for base-disk sampling and 0.001469 for surface-area sampling.
The latter's unit coordinates converge, so its limiting tree uses their limits.

The distinction from the paraboloid is the connectivity of height levels. A
paraboloid height level is one circle, with travel distances only of order $r$.
A positive saddle height level has two separate components that are still
order $r^2$ apart geodesically. The limiting tree is intrinsically
one-dimensional, but representing its branch distances approximately by
Euclidean distances can require a configuration spanning three dimensions.
Intrinsic dimension of the limiting tree and linear span of an MDS embedding
are different concepts.

\subsection{An antipodal contrast}
For the saddle boundary pair $(\pm r,0,r^2)$, reflection in $y=0$ fixes both
endpoints. Uniqueness of the minimizing geodesic forces it onto that meridian.
Its exact distance is
\begin{equation}\label{eq:saddle-antipodal}
 2\int_0^r\sqrt{1+4x^2}\,dx
 =r\sqrt{1+4r^2}+\tfrac12\asinh(2r)\sim2r^2.
\end{equation}
The corresponding paraboloid geodesic is asymptotic to $\pi r$. Conversely,
saddle antipodes on the ruling $x=y$ or $x=-y$ lie on a straight line in
three-dimensional space and have exact distance $2r$. Thus the direction
of a saddle antipodal pair matters. These formulas concern input geodesics;
they do not give a universal fitted MDS distance for an individual saddle pair.

\subsection{A nonplanarity theorem for classical MDS}
\begin{theorem}[Three positive directions in the tree limit]\label{thm:tree-rank}
Let $F$ be the $n\times4$ matrix whose $i$th row is $t_i e_{a(i)}^T$, where
$a(i)$ is the branch index and $e_a$ is a standard basis vector. If $JF$ has
rank four, the centered squared tree-distance matrix has exactly three
positive eigenvalues and one negative eigenvalue. Classical three-dimensional
MDS of the saddle therefore has three nonzero limiting singular values after
division of coordinates by $r^2$ and cannot collapse to a plane or a line.
\end{theorem}
\begin{proof}
Since $t=F\one_4$, expansion of \eqref{eq:tree} gives
\begin{equation}
 T_{ij}^2=t_i^2+t_j^2+2t_it_j-4t_it_j\,1_{\{a(i)=a(j)\}}.
\end{equation}
Double-centering removes the first two terms, yielding
\begin{equation}\label{eq:tree-gram}
 B_T=-\tfrac12JT^{\circ2}J
     =JF(2I_4-\one_4\one_4^T)F^TJ.
\end{equation}
The middle matrix has eigenvalue $2$ on the three-dimensional subspace
orthogonal to $\one_4$ and eigenvalue $-2$ along $\one_4$. If $JF$ has full
column rank, Sylvester's law of inertia transfers these three positive and
one negative directions to $B_T$, with zeros on the remaining subspace.
By Proposition~\ref{prop:tree-bound}, $B(r)/r^4\to B_T$. Classical MDS retains
its three positive directions; their eigenvalues stay positive in the limit.
\end{proof}

The rank condition holds for generic samples covering all four branches with
enough distinct points and is verified for both limiting samples used here.
Degenerate samples or missing branches require separate consideration. For
the present base-disk sample, the limiting classical ratios are
$s_2/s_1=0.918394$ and $s_3/s_2=0.909343$. For the limiting surface-area sample,
they are 0.927947 and 0.923701.

\subsection{What the raw-stress equation proves, and what remains numerical}
With $Y=r^2V$, the exact rescaling is
\begin{equation}\label{eq:tree-stress}
 r^{-4}S_r(r^2V)=\sum_{i<j}
 \left(\norm{V_i-V_j}-\frac{\Delta_{ij}(r)}{r^2}\right)^2
 \longrightarrow\sum_{i<j}(\norm{V_i-V_j}-T_{ij})^2.
\end{equation}
For a fixed finite sample this convergence is uniform on bounded
configurations. Centered global minimizers are bounded: comparing with the
zero configuration bounds stress, which bounds every fitted pair distance;
centering then bounds every coordinate. Any accumulation point of scaled
global minimizers therefore minimizes raw stress for the limiting tree.

This identifies the limiting optimization problem. It does not by itself
prove that every global three-dimensional minimizer of that problem is
nonplanar. The inertia argument in Theorem~\ref{thm:tree-rank} applies to
classical MDS, not automatically to raw stress. The six-start fits supply
numerical evidence of raw-stress nonplanarity for the samples studied here.
Under the theorem's rank condition the tree matrix is not Euclidean, because
its centered squared-distance matrix has a negative eigenvalue. Its minimum
raw stress is therefore positive in every fixed Euclidean output dimension.

\section{Why the original saddle can look like an interval}\label{sec:scaling}

At every finite $r$, the original graph is a saddle extending over the entire
disk. The narrowing in an equal-axis display comes from its changing aspect
ratio. Its horizontal coordinates range over $[-r,r]$, whereas height ranges
over $[-r^2,r^2]$. At $r=32$, the horizontal width is 64 and the vertical span
is 2048, a ratio of 32.

The original experiment's snapshots divide all coordinates by RMS input
geodesic distance, which is of order $r^2$. Uniform division by $r^2$ gives
\begin{equation}
 (x/r^2,y/r^2,z/r^2)=(a/r,b/r,a^2-b^2),
 \qquad a^2+b^2\le1.
\end{equation}
The horizontal coordinates then lie in $[-1/r,1/r]$, while the vertical
coordinate still spans $[-1,1]$. Thus the entire displayed surface is close
to a vertical interval in the ambient coordinate metric.

To retain the familiar appearance, rescale the axes differently:
\begin{equation}
 X=x/r,\quad Y=y/r,\quad Z=z/r^2
 \quad\Longrightarrow\quad
 Z=X^2-Y^2,\quad X^2+Y^2\le1.
\end{equation}
This shows the same saddle-shaped graph for every $r$, but it changes angles
and relative lengths by scaling horizontal and vertical directions differently.
Figure~\ref{fig:scaling} shows the full analytic surface and its boundary under
both choices. The earlier snapshots were finite point clouds, rather than
filled surface meshes; neither display choice removes any part of the domain.

Ambient closeness does not imply geodesic closeness. The points
$(\pm r,0,r^2)/r^2$ approach the same ambient point $(0,0,1)$, yet their
geodesic distance divided by $r^2$ approaches $2$. This is why the original
surface can look like an interval while its geodesic MDS embedding keeps
the two upper branches apart. The two limits concern different metrics.

\section{Discussion and answers to the original questions}\label{sec:discussion}

The raw-stress intuition is correct when it refers to equal relative errors:
a long target then contributes proportionally to the square of its length.
It is incorrect if interpreted as a larger penalty for an equal absolute
error or as a universal tendency to flatten curved inputs. Both the exact
ambient reconstruction and the geodesic saddle experiment rule out that
general conclusion.

For the paraboloid with smooth geodesic input, the leading normalized metric
forgets angle and keeps squared base radius. The stress comparison with a
height-line candidate turns this geometric fact into a proof of the MDS line
limit. The finite-radius shapes are shallower than the original in the
quadratic diagnostic, but the observed coefficients do not trend to zero.
Their infinite-radius limits have not been determined here.

For the saddle, disconnected same-height components survive as different
branches of a tree. The corresponding classical-MDS matrix has three positive
directions, so it retains three-dimensional span. The raw-stress fits show the
same behavior, but a theorem characterizing all global raw-stress minimizers
of the tree remains open in this report.

Boundary-antipodal questions illustrate why the scale and sample assumptions
matter. On the paraboloid, their target separation is $\pi r+O(r^{-1})$, even
though it disappears in the $r^2$ line limit. Global raw stress preserves its
leading coefficient for a fixed generic interior sample plus the pair.
Classical three-dimensional MDS instead has a coefficient determined by a
transverse eigenproblem. The finite sample and population-operator estimates
are different, and neither should be substituted for the fixed-sample raw-stress
coefficient.

The main numerical limitation is that these are coupled single-seed
experiments, not a study of sampling variability. Optimizer convergence is
reported but is not a certificate of global optimality. Smooth-geodesic
validation is stringent numerical evidence, not interval arithmetic. The
most useful further theory would characterize raw-stress minimizers of the
four-branch metric and control joint population/large-radius limits for the
paraboloid's subleading eigenproblem.

\appendix
\section{Smooth-geodesic computation and its justification}\label{app:geodesics}

\subsection{Paraboloid metric and exact primitives}\label{app:parab-geodesics}
In polar coordinates the paraboloid metric is
\begin{equation}\label{eq:parab-metric}
 ds^2=(1+4\rho^2)d\rho^2+\rho^2d\theta^2.
\end{equation}
Along a unit-speed geodesic, Clairaut's first integral is
$c=\rho^2\,d\theta/ds$, with turning radius $\rho=c$ after choosing the
angular orientation so $c\ge0$. The angular equation is
\begin{equation}
 \frac{d\theta}{d\rho}=
 \frac{c\sqrt{1+4\rho^2}}{\rho\sqrt{\rho^2-c^2}}.
\end{equation}
The angular equation, a corresponding elementary primitive, and the Gaussian
curvature are given by \citet[pp.~5--6, Eqs.~(9.2)--(9.4)]{wheeler2016}.
The following equivalent angular expression and the length primitive are
the forms used in the numerical solver. Put
$t=\sqrt{\rho^2-c^2}$, $A=1+4c^2$, and
$v=\asinh(2t/\sqrt A)$. Then
\begin{align}\label{eq:primitives}
 H(c,t)&=2cv+\operatorname{atan2}\!\left(t,c\sqrt{A+4t^2}\right),\\
 L(c,t)&=\frac12t\sqrt{A+4t^2}+\frac14Av.
\end{align}
They measure angle and length from the turning point. In particular,
$dL/dt=\sqrt{1+4c^2+4t^2}$. Pole and zero-angle limits are handled separately
to avoid ambiguity in the angular expression at $c=0$.

Sort endpoint radii as $a\le b$, and let $\alpha\in[0,\pi]$ be the smaller
angular separation. Define $\alpha_0=H(a,\sqrt{b^2-a^2})$. If
$\alpha\le\alpha_0$, solve the monotone-radial branch
\begin{equation}\label{eq:monotone}
 H(c,\sqrt{b^2-c^2})-H(c,\sqrt{a^2-c^2})=\alpha.
\end{equation}
Otherwise solve the turning branch, replacing the difference by a sum. The
length is the corresponding difference or sum of $L$. The implementation uses
62 vectorized bisection steps and the parameterization $c=a\sin\phi$,
$t_a=a\cos\phi$ for numerical stability near a turning point.

\subsection{Why these branches give shortest paths}
A path whose total angular variation exceeds the smaller endpoint separation
can be shortened by compressing its angular coordinate while keeping its
radial profile. In radial arclength coordinates $s_\rho$, the metric is
$ds_\rho^2+\rho(s_\rho)^2d\theta^2$, and the geodesic equation gives
$s_\rho''=\rho\rho'\theta'^2\ge0$. Thus a geodesic has no radial maximum and
at most one radial minimum. It stays within the larger endpoint radius and
has either the monotone or turning form above.

The monotone-branch angular difference increases with $c$. On the turning
branch, each endpoint contribution satisfies
\begin{equation}
 \frac{d}{dc}H(c,\sqrt{\rho^2-c^2})
 =2\asinh\!\frac{2\sqrt{\rho^2-c^2}}{\sqrt{1+4c^2}}
  -\frac{\sqrt{1+4\rho^2}}{\sqrt{\rho^2-c^2}}.
\end{equation}
The derivative of this expression is negative for $0<c<\rho$: the first
term's argument decreases, and the negative second term becomes more negative.
The turning angular sum is therefore strictly concave, starts at $\pi$ as
$c\downarrow0$, and ends at $\alpha_0$ at $c=a$. For $\alpha<\pi$ in the
turning range it crosses the target exactly once. At $\alpha=\pi$, a positive
root, when present, is shorter than the through-apex meridian. Indeed, along
the turning family $dL_{\mathrm{total}}/dc=c\,dH_{\mathrm{total}}/dc$;
integrating by parts between zero and the positive root shows that the length
difference is negative because the intervening angular sum exceeds $\pi$.
Otherwise the meridian is selected.

These arguments select minimizing paths rather than arbitrary winding
geodesics. They, the length primitive, and the MDS asymptotics are supplied
derivations, rather than claims attributed to Wheeler's discussion of general
geodesics on the surface.

\subsection{Saddle shooting and continuation}
The saddle solver integrates \eqref{eq:saddle-ode} in unit base coordinates.
For $P=(ra,rb,r^2(a^2-b^2))$, the equations for $a,b$ contain $4r^2$ in place
of 4. An adaptive Dormand--Prince 5(4) integrator carries position and velocity
together with analytic sensitivities to the two initial velocity components.
Damped Newton shooting solves the two endpoint constraints. The default
relative and absolute integration tolerances are $2\times10^{-10}$ and
$2\times10^{-12}$; the internal endpoint tolerance is $2\times10^{-9}$ in
unit base coordinates.

If the initial normalized velocity is $(v,w)$ at $(a,b)$ and the affine time
interval is $[0,1]$, conservation of metric speed gives the path length as
\begin{equation}
 r\sqrt{v^2+w^2+4r^2(av-bw)^2}.
\end{equation}
Velocities are continued from the preceding radius. Failed solves are retried
with finer continuation steps in radius and endpoints; unresolved pairs stop
the experiment. No graph or mesh distances enter this process. The uniqueness
and convexity argument in Section~\ref{sec:saddle} supplies its global
shortest-path justification once the endpoint problem is accurately solved.

\section{The ambient-distance paraboloid comparison}\label{app:ambient}

The initial comparison used ambient Euclidean dissimilarities rather than the
requested geodesics. It is retained because it isolates two distinct effects:
exact finite-radius reconstruction and apparent narrowing under common scale
normalization. Figures~\ref{fig:ambient-snapshots}--\ref{fig:ambient-sampling}
contain all three figures from that comparison.

\subsection{Exact reconstruction in three dimensions}
For $P_i(r)=(ru_i,r^2q_i)$ with $q_i=\norm{u_i}^2$, set
$\delta_{ij}=\norm{P_i-P_j}$. Choosing $Y_i=P_i$ gives zero raw stress, hence a
global optimum. Every zero-stress fit has the same centered Gram matrix
\begin{equation}
 -\tfrac12J\delta^{\circ2}J=P_cP_c^T
\end{equation}
and is congruent to the original sample, up to translation, rotation, and
reflection. Classical three-dimensional MDS recovers that Gram matrix exactly.
The maximum relative coordinate recovery discrepancy over all twenty cases
was less than $4.83\times10^{-15}$.

For a direct flattening candidate, keep horizontal coordinates fixed and
replace height by $a r^2q_i$. Its pair distances become
\begin{equation}
 \sqrt{r^2\norm{u_i-u_j}^2+a^2r^4(q_i-q_j)^2}.
\end{equation}
For any pair of unequal squared radii, decreasing $|a|$ below one creates a
nonzero residual. The optimum remains $|a|=1$, with the sign representing
reflection. Two-dimensional MDS, by contrast, is planar by construction and
cannot exactly realize the distances of a generic affinely three-dimensional
sample.

\subsection{The normalized limit and population covariance}
The scaled chord distances satisfy
\begin{equation}
 \frac{\delta_{ij}}{r^2}
 =\sqrt{(q_i-q_j)^2+\frac{\norm{u_i-u_j}^2}{r^2}},\qquad
 0\le\frac{\delta_{ij}}{r^2}-|q_i-q_j|\le\frac2r.
\end{equation}
The proof of Theorem~\ref{thm:line} applies with candidate stress bound
$4M/r^2$ instead of $M(1+\pi)^2/r^2$. Thus even exact three-dimensional
reconstruction approaches a line after scale normalization. This is not a
change in the original quadratic coefficient.

For the uniform base-disk population, $q$ is uniform on $[0,1]$ and symmetry
removes the cross-covariances. Therefore
\begin{equation}
 \Cov(P)=\diag\!\left(\frac{r^2}4,\frac{r^2}4,\frac{r^4}{12}\right).
\end{equation}
The leading variance fraction is
\begin{equation}
 \frac{\max\{3,r^2\}}{6+r^2},
\end{equation}
which equals $r^2/(r^2+6)$ for $r\ge\sqrt3$. Small patches are approximately
planar; large normalized patches are approximately linear. This is not
monotone flattening into a plane.

\subsection{Numerical estimates and absolute versus relative error}
\input{tables/ambient_sweep.tex}
At $r=64$, surface-area sampling gives two-dimensional raw-stress relative
RMSE 0.003688 and leading variance fraction 0.998431. Simply deleting height
and optimally rescaling the original horizontal coordinates gives relative
RMSE 0.6415 for the base-disk sample at that radius. This is one specified
flat-disk candidate, not a lower bound over all planar embeddings.

Between radii 4 and 64, the selected base-disk two-dimensional raw-stress fit
improves in relative RMSE from 0.1141 to 0.002611, while its absolute RMSE grows
from 0.8828 to 4.4392. A small relative error on a rapidly expanding distance
scale does not mean the absolute errors are shrinking.

\section{Validation, provenance, and reproduction}\label{app:reproduction}

\subsection{Checks on the distance solvers and embeddings}
For the paraboloid, 50 validation pairs across five radii were compared with
adaptive quadrature of the original angle and length integrals and with
independent Cartesian geodesic integration. The largest relative length
disagreement was less than $3.05\times10^{-15}$, angular disagreement less
than $1.64\times10^{-14}$ radians, and endpoint error divided by disk radius
less than $3.98\times10^{-11}$. Pole, duplicate, same-meridian, and
opposite-meridian cases were also checked.

For the saddle, 60 validation pairs at five radii were checked using independent
SciPy DOP853 integration and boundary-value collocation. Maximum relative
length disagreement was $7.87\times10^{-9}$ or less; maximum independent
endpoint discrepancy divided by disk radius was $8.11\times10^{-7}$ or less.
The latter is an external validation result and is distinct from the tighter
internal shooting residual. Maximum relative energy drift was less than
$6.32\times10^{-11}$. Exact straight rulings and reflection-axis meridian
integrals provide additional checks.

Every full geodesic matrix was checked against ambient chord lower bounds,
explicit surface-path upper bounds, and all triangle inequalities. The saddle
also passed the tree bounds. The triangle-inequality tolerance was $10^{-8}$
times RMS input distance for the paraboloid and $10^{-7}$ for the saddle.
The default saddle sweep contains 573,600 unordered pair distances across
its twenty cases, and so does the paraboloid sweep.

The analytic raw-stress gradient was independently checked by finite differences
in the ambient experiment, with an $L^2$ discrepancy of $1.82\times10^{-8}$.
The largest selected normalized-objective gradient components were below
$4.66\times10^{-9}$ for ambient two-dimensional fits,
$6.58\times10^{-9}$ for paraboloid geodesic fits, and
$4.95\times10^{-9}$ for saddle geodesic fits. Selected three-dimensional raw
fits were checked to have no larger raw stress than their two-dimensional
and classical three-dimensional counterparts, within rounding tolerance.
These comparisons are checks on the recorded solutions, not proofs that a
global optimizer was found.

\subsection{Reproduction without rerunning the analysis for prose changes}
Result generation and report generation are separate. The report consumes
the existing CSV, NPZ, vector-figure, and manifest artifacts. Its tables are
generated directly from the CSV results. The axis-scaling illustration is a
rendering of the analytic saddle equation and does not rerun MDS.

From the repository root, the complete analysis and report can be reproduced by
the following targets:
\begin{quote}\small\ttfamily
make -C tools/experiments/paraboloid-mmds-radius run\\
make -C tools/experiments/paraboloid-mmds-radius geodesic\\
make -C tools/experiments/paraboloid-mmds-radius antipodal\\
make -C tools/experiments/saddle-mmds-radius run\\
make -C tools/reports/geodesic-mds pdf
\end{quote}
The report build verifies the experiment bundles, checksums its consumed
inputs, compiles LaTeX, and checks reference and citation resolution. It
records a build-generated timestamp in the \texttt{America/New\_York} timezone.
All four external citations have BibTeX metadata and claim-level source
evidence in the accompanying \texttt{citation\_verification.html}.

The recorded numerical environment uses Python 3.9.6, NumPy 1.24.4, SciPy 1.9.1,
Matplotlib 3.7.5, and Numba 0.57.1 for saddle geodesics. Actual versions and
checksums are preserved in the experiment and report manifests. The experiments
do not change the R package or its public behavior.

\clearpage
\section{Complete figure atlas}\label{app:figures}

The following fifteen plates include every distinct figure generated during
the investigation. A PNG and PDF export of the same plot are one figure;
vector versions are used whenever available. The final scaling illustration
has been regenerated as a vector figure with the same analytic content and
less crowded ticks. Larger plate pages preserve the native label sizes.

Figures~\ref{fig:p3d}--\ref{fig:pshape} concern the requested paraboloid
geodesics. Figures~\ref{fig:s3d}--\ref{fig:scompare} concern the saddle and its
comparison with the paraboloid. Figure~\ref{fig:scaling} resolves the axis-scaling
question. The final three figures are the ambient-distance comparison and
must not be interpreted as geodesic results.

\plate{geodesic_3d_snapshots.pdf}{fig:p3d}{Paraboloid with smooth geodesic input:
original coordinates, classical three-dimensional MDS, and raw-stress
three-dimensional MDS at four radii. The sample is uniform in the base disk,
$n=240$. Color denotes squared base radius divided by $r^2$. Each case is
centered, rigidly aligned, and divided by its RMS input geodesic distance;
axes use equal units and common limits. Large-radius narrowing is a relative
scale effect, not loss of the original surface's domain.}

\plate{geodesic_diagnostics.pdf}{fig:pdiag}{Paraboloid geodesic radius sweep:
relative distance RMSE and singular-value ratios for classical 3D, raw-stress
3D, and raw-stress 2D MDS. Definitions are in Section~\ref{sec:definitions}.
The falling $s_2/s_1$ at large radius identifies the line limit; $s_3/s_2$
remaining substantial means the two smaller directions shrink together.
The error curves compare objectives and imposed output dimensions.}

\plate{geodesic_2d_snapshots.pdf}{fig:p2d}{Paraboloid geodesic raw-stress MDS
when two-dimensional output is imposed. Coordinates are shown in principal
component (PC) axes and divided by RMS input geodesic distance, with equal
spatial units. Color denotes squared normalized base radius. Radial ordering
becomes the dominant separation at large radius.}

\plate{geodesic_sampling.pdf}{fig:psampling}{Paraboloid geodesic raw-stress
3D MDS under uniform base-disk and uniform surface-area sampling. Each curve
uses coupled quantiles across radius. Both measures show decreasing relative
distance error and increasing concentration toward a line at large radius;
their numerical rates differ. These are single-sample curves, not uncertainty
intervals.}

\plate{geodesic_shape.pdf}{fig:pshape}{Paraboloid shape regression after rigid
alignment without scale fitting: $z=a(x^2+y^2)+bx+cy+d$. The original coefficient
is one. Classical and raw-stress MDS give shallower fitted shapes, but the
coefficients do not tend toward zero over the tested radii. The right panel
uses a focused $R^2$ scale beginning at 0.9. The regression describes a point
cloud; it does not identify an intrinsic curvature.}

\plate{saddle_3d_snapshots.pdf}{fig:s3d}{Saddle with smooth geodesic input:
original coordinates, classical 3D MDS, and raw-stress 3D MDS. Colors identify
the four branches defined in Section~\ref{sec:saddle}. Each case is centered,
rigidly aligned, and divided by its RMS geodesic distance, with equal spatial
units and common limits. The original ambient display narrows, while MDS
separates geodesically distant branches in three dimensions.}

\plate{saddle_diagnostics.pdf}{fig:sdiag}{Saddle geodesic radius sweep.
Relative distance RMSE compares classical 3D, raw-stress 3D, and raw-stress 2D
fits. Neither $s_2/s_1$ nor $s_3/s_2$ tends toward zero for the 3D fits over
the large-radius range. The input approaches a non-Euclidean tree metric,
so a positive limiting relative fitting error is compatible with the theory.}

\plate{saddle_2d_snapshots.pdf}{fig:s2d}{Saddle geodesic raw-stress MDS with
two-dimensional output imposed. The coordinates use principal-component
axes, equal units, and RMS input-distance normalization. Four narrow branches
develop in the plane. Their planar appearance is an output constraint;
the corresponding three-dimensional fit retains an additional substantial
direction and achieves lower stress.}

\plate{saddle_sampling.pdf}{fig:ssampling}{Saddle geodesic raw-stress 3D MDS
under both sampling measures. Relative distance error approaches different
positive levels, while the third-to-second singular-value ratio remains near
one. Nonplanarity is therefore not an artifact of choosing uniform base-disk
sampling rather than uniform surface area.}

\plate{saddle_shape.pdf}{fig:sshape}{General saddle quadratic fit
$z=ax^2+bxy+cy^2+dx+ey+f$ after rigid alignment. The coefficient norm is
$\sqrt{a^2+c^2+b^2/2}$, with original value $\sqrt2$. Its decrease does not
imply a plane: coefficients carry inverse-length units, while the embedding
expands. The accompanying $R^2$ measures only the descriptive quadratic fit.}

\plate{saddle_vs_paraboloid.pdf}{fig:scompare}{Direct comparison of raw-stress
3D MDS using smooth geodesic distances and the same 240 base-disk samples.
The paraboloid's $s_2/s_1$ falls toward zero, whereas the saddle's remains large.
The right panel alone would not distinguish these outcomes because two
shrinking transverse axes of a line-like cloud can remain comparable. Both
ratios are needed to interpret normalized shape.}

\plate{saddle_axis_scaling.pdf}{fig:scaling}{The same full analytic saddle at
$r=32$, including its disk boundary. Left: all coordinates are divided by
$r^2$, preserving their relative units and producing a narrow ambient view.
Right: horizontal coordinates are divided by $r$ and height by $r^2$,
recovering the familiar graph $Z=X^2-Y^2$ on the unit disk. The right-hand
scaling changes relative lengths; it is useful for displaying the surface,
but is not the metric normalization used for MDS comparisons.}

\plate{embedding_snapshots.pdf}{fig:ambient-snapshots}{Ambient-distance
comparison only: exact 3D reconstruction and raw-stress 2D MDS of the paraboloid.
The inputs here are Euclidean chords, not smooth geodesics. Colors identify
squared normalized base radius. Common RMS input-distance scaling makes the
exactly reconstructed 3D cloud look narrow at large $r$, even though it is
congruent to the original sample in its original units.}

\plate{radius_diagnostics.pdf}{fig:ambient-diag}{Ambient-distance comparison
only: relative distance RMSE for 2D methods and a height-line candidate,
alongside leading principal-component variance fractions. The population
curve is $\max\{3,r^2\}/(6+r^2)$ under uniform base-disk sampling. The exact
3D reconstruction has zero stress to rounding and still approaches a line
after scale normalization.}

\plate{sampling_comparison.pdf}{fig:ambient-sampling}{Ambient-distance
comparison only: the radius dependence of two-dimensional raw-stress MDS
under uniform base-disk and uniform surface-area sampling. The two measures
give the same qualitative large-radius line concentration. Relative error
is normalized by the RMS chord distance, so it need not track absolute
distance error.}

\clearpage
\bibliographystyle{plainnat}
\bibliography{references}
\end{document}
